\documentclass{article}

\title{computer networks}
\author{alexander}
\date{\today}

\begin{document}

\maketitle

I am writing this document as a passion project for myself. This document will be writing mostly from my memory with some help from LLM(s).\\

A switch, for the most part, handles layer two traffic. Switches use the Address Resolution Protocol(ARP) to determine where on a Local Area Network(LAN) a given device with a certain Media Access Control(MAC) address may be. The ARP sends what is known as broadcasts out of all of its ports and waits for a reply from the machine with a given MAC address. Switches can act as a physical segmentation of a LAN and they can implemnet something known as a VLAN to logically segment that same network. The ARP can create broadcast storms and other issues locally so there is another useful protocol that switches use known as spanning tree protocols. There are a variety of spanning tree protocols: STP, RSTP, MSTP, PSTP, ect. each with their own nuiances and benifits. There are also a variety of security protocols switches implament because after all switches are like the first layer into a network making them a valuable asset for attacks to target.

Routers can do many more things than switches in my opinion. Routers will look at layer two information like a switch, but that really only gets you to/from the next hop. Routers will look at layer three information. When I say layer three, I am referring to the OSI model (IP layer). There are two types of IP: IPv4, IPv6 each with their own reasons for exisiting. Routers can maintain routing table which are similar in some ways to the MAC address tables that are maintaed by switches. Routing tables can be manually populated with static routes set up by an admin. Routers can also alter the routing table dynamically with routing protocols. There are many routing protocols with their own caveats: OSPF, IEGRP, IS-IS, BGP, RIP, ect. Routers can have some what of a secuirity feature called Access Control Lists(ACLs). These list can filter traffic a little bit by looking at different encapsulation layers.

Do not think of "internet speed" as one number. Think of a network path as a system with a capacity, delay, queues, loss, and variability. Bandwidth is the capacity of a link or path to carry data (bit/s, Mbps, Gbps). A rural area may have a bandwidth of 50Mbps because of the infrastructure (DSL). In a larger metroplex or near a university the capacity can be exceeding 1Gbps. Throughput is the actual rate of successful data transfer at a particular time under particular conditions while bandwidth is the what the link is capable of providing. Delay is the amount of time it takes for a packet to experience some part of its journey through the network. This is a broad term....A packet can experience several kinds of delay. Delay is the time introduced while the packet is being handled and transported. Latency is the term you will encounter much more often when talking about the user experience. Latency is how long communication takes. You ping a server the Round Trip Time (RTT) is 20ms.

processing delay: device processing the packet
transmission delay: time required to put the packet's bits onto the outgoing link
propagation delay: time for the signal to physically travel through the medium
queing delay: time the packet spends waiting in a queue

delay is the general networking concept...latency is the end-to-end responsiveness you experience/measured between endpoints.

Jitter -> latency/delay variation










\end{document}
